\scanpage{6}\mspage{2}
l'élément $x\wedge y$ de $(\det M)^*$ ne dépend que de la permutation circulaire $\omega\in\underline{\omega}_A$ définie par l'ordre total choisi sur $A$, on trouve donc une application canonique
\[
\underline{\omega}_A\xrightarrow{\alpha_A}(\det M)^*\tag*{(10)}
\]
satisfaisant
\[
\alpha_A(\omega^{-1})=-\alpha_A(\omega)\tag*{(11)}
\]
de sorte que $\alpha_A$ peut aussi s'interpréter comme un élément
\[
\alpha_A\in\left(\det_{\Lambda}(M)\wedge_{\mu_2}\underline{\omega}_A\right)^*\simeq
\det_{\Lambda}(M)^*\wedge_{\mu_2}\underline{\omega}_A.\tag*{(12)}
\]
\deleted{On} Si $M$ est ``jaugé'' i.e. muni d'une \uline{base} de $\det M$, soit $e_M$, on pose
\[
\left\{\begin{aligned}
S\widetilde F(M)&=\left\{A\in\widetilde F(M)\mid\alpha_A\in
\underbrace{\{\pm e_M\}\wedge_{\mu_2}\underline{\omega}_A}_{\simeq\,\underline{\omega}_A}\right\}\\
\deleted{S\widetilde F^+(M)}&\deleted{=\{A\in\widetilde F(M)\mid\alpha_A\unsure{\in}\ \in}
\end{aligned}\right.\tag*{(13)}
\]
i.e. $S\widetilde F(M)$ est formé des $A\in\widetilde F(M)$ tels que l'image de $\underline{\omega}_A$ par (\deleted{\unsure{4|1}}10) soit $\{\pm e_M\}$. Si $2\ne0$ dans $\Lambda$, i.e si $\Lambda$ n'est pas une $\mathbb F_2$-algèbre, alors \deleted{(10) est injectif, \unsure{et|donc|dès}}
\[
\alpha_A:\ \underline{\omega}_A\simeq\{\pm e_M\}
\]
est une bijection, ce qui implique que
% archive sheet number: 5
